% BHU PYQ -- Manually transcribed & error-corrected
% Paper: BSC-02A: Elements of Geology - I (Ancillary)
% Exam: B.Sc. (Hons.) Semester II Examination, 2013-14

\documentclass[11pt,a4paper]{article}
\usepackage[a4paper,top=2.5cm,bottom=2.5cm,left=2.8cm,right=2.8cm,headheight=14pt]{geometry}
\usepackage{amsmath,amssymb}
\usepackage[T1]{fontenc}
\usepackage[utf8]{inputenc}
\usepackage{lmodern,microtype,fancyhdr,enumitem,hyperref,lastpage,parskip}

\hypersetup{colorlinks=true,urlcolor=black,linkcolor=black,
  pdftitle={BSC-02A Elements of Geology - I (Ancillary) -- BHU B.Sc. Sem II 2013-14}}
\pagestyle{fancy}\fancyhf{}
\renewcommand{\headrulewidth}{0.4pt}\renewcommand{\footrulewidth}{0.4pt}
\fancyhead[L]{\small   \textit{Banaras Hindu University}}
\fancyhead[R]{\small   \textit{BSC-02A: Elements of Geology - I (Ancillary)}}
\fancyfoot[C]{\small Page  \thepage\ of \pageref{LastPage}}


\newlist{parts}{enumerate}{2}
\setlist[parts,1]{label=(\alph*),leftmargin=*,itemsep=5pt,topsep=3pt}

\newcommand{\pts}[1]{\hfill   \textit{[#1]}}


\begin{document}

\begin{center}
  {\large\bfseries BANARAS HINDU UNIVERSITY}\\[2pt]
  {\normalsize B.Sc. (Hons.) --- Semester II Examination, 2013-14}\\[6pt]
  \rule{0.8\linewidth}{0.5pt}\\[6pt]
  {\large\bfseries Paper: BSC-02A --- Elements of Geology - I (Ancillary)}\\[4pt]
  {\normalsize Time: 3 Hours \hfill Maximum Marks: 70}\\[2pt]
  {\small Ref. No.: BS/Sem II/17}
\end{center}
\rule{\linewidth}{0.4pt}
\smallskip



\noindent   \textit{Instructions: Answer   \textbf{five} questions in all, selecting at least   \textbf{two} from each of Sections A and B. Section-C is compulsory. All questions carry equal marks.}
\bigskip

\bigskip


\noindent {\large\bfseries SECTION --- A}
\rule{\linewidth}{0.4pt}\smallskip




\noindent   \textbf{Question 1.} Describe the geological work of a river. \pts{14}

\medskip


\noindent   \textbf{Question 2.} Write a detailed note on   \textbf{any two} of the following: \pts{$2  \times 7 = 14$}

\begin{parts}
  \item Parts of folds
  \item Ice ages
  \item Dip and strike
\end{parts}

\medskip


\noindent   \textbf{Question 3.} Discuss the relevance and importance of Earth System Science vis-à-vis its different branches. \pts{14}

\medskip


\noindent   \textbf{Question 4.} What are sedimentary rocks? Write a detailed note on their classification giving suitable examples. \pts{14}

\bigskip


\noindent {\large\bfseries SECTION --- B}
\rule{\linewidth}{0.4pt}\smallskip




\noindent   \textbf{Question 5.} Explain the mechanism of an earthquake with a neat sketch. Discuss the societal and environmental impact of earthquakes. \pts{14}

\medskip


\noindent   \textbf{Question 6.} Write short notes on   \textbf{any four} of the following: \pts{$4  \times 3\frac{1}{2} = 14$}

\begin{parts}
  \item Barchans
  \item Playa
  \item Chemical weathering
  \item Moraines
  \item Nunataks
  \item Exfoliation
\end{parts}

\medskip


\noindent   \textbf{Question 7.} How do you distinguish between   \textbf{any two} of the following mineral pairs based on their diagnostic physical properties? \pts{$2  \times 7 = 14$}

\begin{parts}
  \item Quartz and Olivine
  \item Garnet and Muscovite
  \item Biotite and Tourmaline
  \item Plagioclase and Hornblende
\end{parts}

\bigskip


\noindent {\large\bfseries SECTION --- C}
\rule{\linewidth}{0.4pt}\smallskip




\noindent   \textbf{Question 8.} Explain the following in 2 to 3 sentences: \pts{$14  \times 1 = 14$}

\begin{parts}
  \item Gabbro
  \item Asthenosphere
  \item Plate tectonics
  \item Normal fault
  \item Meteorology
  \item Lustre in minerals
  \item Mohs' scale
  \item Richter scale
  \item Continental drift
  \item Thrust fault
  \item Mass extinction
  \item Denudation
  \item ``Present is the key to the past''
  \item Marble
\end{parts}

\vfill\rule{\linewidth}{0.4pt}

\begin{center}\small
\href{https://bhu-exam.vercel.app/}{Click here to visit our website}\\[4pt]
\href{https://me-aryan.vercel.app/}{Click here to visit our contributor}\\[4pt]
\href{https://quashan-me.vercel.app/}{Click here to visit our contributor}
\end{center}
\end{document}
